% ============================================================
% GSPH-FC — Springer LNCS | WEA 2026
% Compilar: pdflatex -> bibtex -> pdflatex -> pdflatex
% ============================================================
\documentclass[runningheads]{llncs}
\usepackage[T1]{fontenc}
\usepackage[utf8]{inputenc}
\usepackage{graphicx}
\usepackage{amsmath,amssymb}
\usepackage{booktabs}
\usepackage{algorithm}
\usepackage[noend]{algpseudocode}
\usepackage{url}

\begin{document}

% ---- Título ----
\title{GSPH-FC: A Geometric Space Partitioning Heuristic
with Frontier Connectivity for the Euclidean TSP}
\titlerunning{GSPH-FC: Geometric Space Partitioning for the TSP}

% ---- Autores ----
\author{%
  Martín Gómez Navarro\inst{1} \and
  Gustavo Esteban Gatica\inst{1} \and
  Rodrigo Linfati\inst{2} \and
  Juan Carlos Figueroa-García\inst{3} \and
  Carlos Alberto Franco Franco\inst{3} \and
  Daniel Morillo-Torres\inst{4}}
\authorrunning{M. Gómez Navarro et al.}
\institute{%
  Faculty of Engineering, Universidad Andres Bello, Santiago, Chile\\
  \email{m.gmeznavarro@uandresbello.edu}\\
  \email{ggatica@uandresbello.edu}
  \and
  Faculty of Engineering, Universidad del Bío-Bío,
  Concepción, Chile\\
  \email{rlinfati@ubiobio.cl}
  \and
  Faculty of Engineering, Universidad Distrital
  Francisco José de Caldas, Bogotá, Colombia\\
  \email{jcfigueroag@udistrital.edu.co}\\
  \email{cafrancof@udistrital.edu.co}
  \and
  Faculty of Engineering, Pontificia Universidad Javeriana de Cali,
  Cali, Colombia\\
  \email{daniel.morillo@javerianacali.edu.co}}

\maketitle

% ================================================================
\begin{abstract}
We propose \textbf{GSPH-FC} (\textit{Geometric Space Partitioning
with Frontier Connectivity}), a hierarchical heuristic for the
Euclidean \textit{Traveling Salesman Problem} (TSP). The method
decomposes the instance into four quadrants via the bounding-box
midpoint, builds and refines local routes within each quadrant
using Nearest Neighbour with 2-opt, and assembles a global tour
through minimum-cost frontier connections. An optional polishing
phase applies 2-opt with Don't-Look Bits, Or-opt chain relocation,
and a lightweight ILS with double-bridge perturbations under a
wall-clock budget. We report a systematic evaluation on 55 TSPLIB
instances ($n$ up to 5\,934) under \texttt{EUC\_2D}, achieving a
typical GAP of 3--12\% over the Best Known Solution with sub-second
times for $n\!\le\!300$, near-constant peak memory ($\approx$4.75\,MB)
across instance sizes, and natural 4-thread parallelism reducing the
effective 2-opt cost to $O(n^2/4)$.
\keywords{2-opt \and Frontier connectivity \and Geometric decomposition
\and Space partitioning \and TSPLIB benchmarks \and
Traveling Salesman Problem}
\end{abstract}

% ================================================================
\section{Introduction}
% ================================================================

The Euclidean \textit{Traveling Salesman Problem} (TSP) is a
cornerstone of combinatorial optimization: find a minimum-length
Hamiltonian cycle over $n$ points in the plane~\cite{applegate2006tsp,laporte1992tsp}.
While exact methods guarantee optimality, their exponential
complexity makes them impractical beyond a few hundred
nodes~\cite{garey1979computers}; heuristics and metaheuristics
therefore remain the dominant tools for large-scale instances,
trading optimality for speed and scalability~\cite{mardones2023}.

We introduce \textbf{GSPH-FC} (\textit{Geometric Space Partitioning
with Frontier Connectivity}), a hierarchical heuristic that exploits
the spatial structure of Euclidean instances~\cite{karp1977probabilistic}. The algorithm operates
in \textbf{three main stages}: (i)~spatial decomposition into four
quadrants via the bounding-box midpoint, (ii)~independent Nearest
Neighbour + 2-opt per quadrant, and (iii)~global assembly by
minimum-cost frontier pairs along quadrant boundaries. This
divide-and-conquer structure reduces the effective search cost,
enables natural parallelism, and yields low-crossing tours.

\paragraph{Contributions.}
\begin{enumerate}
  \item \textbf{GSPH-FC algorithm:} geometric design and full
    pseudocode for all subroutines, with $O(n^2)$ complexity
    analysis.
  \item \textbf{Systematic benchmarking:} evaluation on 55 TSPLIB
    instances ($n$ up to 5\,934) reporting GAP\%, wall-clock time,
    and peak memory profiling.
  \item \textbf{Critical analysis:} SWOT assessment and directions
    for multi-level QML-GSPH.
\end{enumerate}

The remainder is organized as follows. Section~\ref{sec:related}
reviews related work. Section~\ref{sec:formulation} presents the
formulation. Section~\ref{sec:algorithm} details pseudocode.
Section~\ref{sec:experiments} reports results.
Section~\ref{sec:analysis} provides critical analysis, and
Section~\ref{sec:conclusions} concludes.

% ================================================================
\section{Related Work}
\label{sec:related}
% ================================================================

We survey the algorithmic families that form the background of GSPH-FC,
covering constructive and improvement heuristics as well as
population-based and trajectory metaheuristics.

\subsection{Constructive and Improvement Heuristics}

Nearest Neighbour (NN)~\cite{kizilatecs2013nearest} is the simplest
constructive method: starting from an arbitrary city, it greedily
selects the closest unvisited city at each step ($O(n^2)$). Despite
practical speed, NN solutions are typically 20--25\% above optimal
and quality depends strongly on the seed. The 2-opt
heuristic~\cite{croes1958method} improves a given tour by reversing
path segments that reduce total length, serving as the standard
local-search baseline and commonly seeded with NN solutions;
Johnson and McGeoch~\cite{johnson1997tsp} provide an extensive
empirical study of its behavior on TSPLIB instances.
The Christofides--Serdyuk algorithm~\cite{christofides1976worst}
achieves a classical $3/2$ approximation guarantee for metric TSP.
The Lin-Kernighan heuristic~\cite{lin1973effective} generalizes
2-opt and 3-opt through an adaptive sequential search; Helsgaun's
LKH~\cite{helsgaun2000effective,helsgaun2017extension} further
incorporates Lagrangian penalties and advanced candidate structures,
reaching near-optimal results.

\subsection{Metaheuristics}

Metaheuristics~\cite{blum2003metaheuristics,sorensen2015metaheuristics}
escape local optima at the cost of longer runtimes. Genetic
Algorithms~\cite{GENETICS,goldberg1985alleleslociand} evolve
populations of tours via crossover operators (OX, PMX) and mutation.
Simulated Annealing~\cite{kirkpatrick1983optimization} accepts
worsening moves with temperature-dependent probability, enabling
escape from local optima. Tabu Search~\cite{glover1989tabu} uses
short-term memory to avoid cycling. Ant Colony
Optimization~\cite{dorigo1992optimization,dorigo1996ant} constructs
solutions probabilistically guided by pheromone trails that
reinforce high-quality routes. While these methods often achieve
smaller gaps than simple heuristics, they require careful parameter
tuning and longer time budgets. Table~\ref{tab:comparison}
summarizes the main algorithmic families.

\begin{table}
\caption{Algorithmic families for the TSP: complexity, scalability,
and solution quality.}\label{tab:comparison}
\centering\small
\begin{tabular}{@{}lllll@{}}
\toprule
\textbf{Method} & \textbf{Type} & \textbf{Complexity}
  & \textbf{Scalability} & \textbf{Quality} \\
\midrule
Brute Force           & Exact         & $O(n!)$           & Very low  & Optimal \\
Branch \& Bound       & Exact         & Exponential       & Low       & Optimal \\
Dynamic Programming   & Exact         & $O(n^2 2^n)$      & Low       & Optimal \\
Mixed Integer Prog.   & Exact         & Exponential       & Low--Med  & Optimal \\
\midrule
Nearest Neighbour     & Constructive  & $O(n^2)$          & High      & Low     \\
2-opt                 & Improvement   & $O(n^2)$/iter     & Medium    & Medium  \\
Lin--Kernighan        & Improvement   & $O(n^{2.2})$      & High      & High    \\
\midrule
Genetic Algorithms    & Metaheuristic & Variable          & High      & Good    \\
Simulated Annealing   & Metaheuristic & Variable          & Medium    & Good    \\
Ant Colony Opt.       & Metaheuristic & Variable          & Med--High & Good    \\
\midrule
\textbf{GSPH-FC (ours)} & Heuristic  & $O(n^2)$ & \textbf{Very high} & \textbf{Good} \\
\bottomrule
\end{tabular}
\end{table}

% ================================================================
\section{Mathematical Formulation}
\label{sec:formulation}
% ================================================================

Let $G=(V,E)$ be a complete undirected graph where
$V=\{v_1,\dots,v_n\}\subset\mathbb{R}^2$ with coordinates
$(x_i,y_i)$. Edge costs follow the TSPLIB \texttt{EUC\_2D}
convention~\cite{reinelt1991tsplib}:
\begin{equation}
d(v_i,v_j)=\left\lfloor\sqrt{(x_i-x_j)^2+(y_i-y_j)^2}+0.5\right\rfloor.
\label{eq:dist}
\end{equation}
The TSP seeks a Hamiltonian cycle
$\pi=(v_{\pi(1)},\dots,v_{\pi(n)},v_{\pi(1)})$ minimizing:
\begin{equation}
\min_{\pi}\;L(\pi)=\sum_{i=1}^{n}
d\!\left(v_{\pi(i)},v_{\pi(i\bmod n+1)}\right).
\label{eq:tsp}
\end{equation}

\paragraph{Scope and assumptions.}
$V\subset\mathbb{R}^2$, $G$ complete undirected, cost $d$ symmetric
Euclidean with \texttt{EUC\_2D} rounding. We seek a Hamiltonian
cycle $\pi$ minimizing $L(\pi)$.

\subsection{Geometric Decomposition}

The bounding box of $V$ yields midpoints:
\begin{equation}
x_m=\tfrac{x_{\min}+x_{\max}}{2}, \quad
y_m=\tfrac{y_{\min}+y_{\max}}{2}.
\end{equation}
$V$ is partitioned into four quadrants:
\begin{align}
Q_1&=\{v\in V: x_v\le x_m,\; y_v>y_m\}  &&\text{(NW)},\\
Q_2&=\{v\in V: x_v>x_m,\; y_v>y_m\}    &&\text{(NE)},\\
Q_3&=\{v\in V: x_v\le x_m,\; y_v\le y_m\}&&\text{(SW)},\\
Q_4&=\{v\in V: x_v>x_m,\; y_v\le y_m\} &&\text{(SE)}.
\end{align}

\subsection{Intra-Quadrant Construction and Refinement}

For each non-empty $Q_k$ with $|Q_k|\ge 2$, a path $P_k$ is built
by NN starting from seed $s\in Q_k$:
\begin{equation}
p_{i+1}=\arg\min_{v\in Q_k\setminus\{p_1,\dots,p_i\}}d(p_i,v).
\end{equation}
Two 2-opt passes are then applied. A swap is accepted if:
\begin{equation}
d(p_i,p_j)+d(p_{i+1},p_{j+1})<d(p_i,p_{i+1})+d(p_j,p_{j+1})
\;\Rightarrow\;\text{reverse }(p_{i+1},\dots,p_j).
\end{equation}

\subsection{Inter-Quadrant Frontier Connectivity}

Adjacent quadrant paths are joined via minimum-cost frontier pairs.
For a vertical boundary at $x_m$ between $Q_a$ and $Q_b$:
\begin{equation}
F_a(\varepsilon)=\{u\in Q_a:|x_u-x_m|<\varepsilon\},\quad
F_b(\varepsilon)=\{v\in Q_b:|x_v-x_m|<\varepsilon\},
\end{equation}
\begin{equation}
(u^*,v^*)=\arg\min_{u\in F_a,\,v\in F_b}d(u,v).
\label{eq:frontier}
\end{equation}
Horizontal boundaries are handled analogously.
$\varepsilon$ is adapted to avoid empty frontier sets.

\subsection{Global Tour Assembly}

Sub-tours are concatenated in clockwise order
$Q_1\!\to\!Q_2\!\to\!Q_4\!\to\!Q_3$. For each consecutive pair
$(R_i,R_{i+1})$, the orientation minimizing the linking cost is
selected:
\begin{equation}
\begin{split}
c(R_i,R_{i+1})=\min\bigl\{
  &\,d(f(R_i),f(R_{i+1})),\;d(f(R_i),\ell(R_{i+1})),\\
  &\,d(\ell(R_i),f(R_{i+1})),\;d(\ell(R_i),\ell(R_{i+1}))\bigr\},
\end{split}
\end{equation}
where $f(\cdot)$ and $\ell(\cdot)$ denote first and last nodes.

\subsection{Advanced Local-Search Polishing}

The optional polishing phase combines three operators under budget
$T_{\text{bgt}}$:

\paragraph{2-opt with Don't-Look Bits (DLB)~\cite{bentley1992fast}.}
For each active node $v_a$ ($\text{dlb}[a]=\mathrm{false}$) and
each candidate $v_c\in C(v_a)$ ($k$ nearest neighbours):
\begin{equation}
\Delta_{2\text{opt}}=
  d(v_a,v_{a+1})+d(v_c,v_{c+1})
 -d(v_a,v_c)-d(v_{a+1},v_{c+1}).
\end{equation}
If $\Delta_{2\text{opt}}>0$ the swap is applied; otherwise
$\text{dlb}[a]\gets\mathrm{true}$.

\paragraph{Or-opt~\cite{or1976traveling}.}
Chains of length $\ell\le L_{\max}$ are relocated from $i$ to $j$:
\begin{equation}
\begin{split}
\Delta_{\text{or}}=\;
  &d(v_{i-1},v_i)+d(v_{i+\ell-1},v_{i+\ell})+d(v_j,v_{j+1})\\
  &-d(v_{i-1},v_{i+\ell})-d(v_j,v_i)-d(v_{i+\ell-1},v_{j+1}).
\end{split}
\end{equation}

\paragraph{ILS Double-Bridge~\cite{lourencco2003ils,martin1991large}.}
Given cut points $0<a<b<c<d<n$:
\begin{equation}
\begin{split}
\pi'\!={}&(\pi_{1:a},\,\pi_{c+1:d},\,\pi_{b+1:c},\,\\
         &\phantom{(}\pi_{a+1:b},\,\pi_{d+1:n}).
\end{split}
\end{equation}
After perturbation, 2-opt DLB is re-applied and the improved
incumbent is retained.

\subsection{Complexity Analysis}

\begin{itemize}
  \item \textbf{Spatial partition:} $O(n)$.
  \item \textbf{NN per quadrant:} $O(n_k^2)$, $n_k\approx n/4$.
    Total: $\sum_k O(n_k^2)\approx O(n^2/4)$.
  \item \textbf{2-opt per quadrant (×2):} $O(n_k^2)$ per pass.
    Highly parallelizable across 4 threads.
  \item \textbf{KNN candidate sets:} $O(n\log n)$ build,
    $O(nk)$ queries.
  \item \textbf{Frontier connectivity:} $O(|F_a||F_b|)$;
    $O(k)$ with cross-frontier $k$-NN index.
  \item \textbf{Polishing (2-opt DLB + Or-opt + ILS):} $O(n^2)$.
\end{itemize}

The dominant term is $T(n)=O(n^2)$, matching state-of-the-art
2-opt-based heuristics, with a practical $\approx 4\times$
reduction from quadrant decomposition and natural 4-thread
parallelism.

\paragraph{Finite convergence.}
Each accepted local move strictly reduces $L(\pi)$; the tour space
is finite, so local-search phases converge in finite time to a
local optimum for the chosen neighborhood.

% ================================================================
\section{Algorithm and Pseudocode}
\label{sec:algorithm}
% ================================================================

Algorithm~\ref{alg:gsph} presents the top-level GSPH-FC procedure.
Algorithms~\ref{alg:nn}--\ref{alg:dlb} detail all subroutines.

\begin{algorithm}
\caption{GSPH-FC — Top-level procedure, $O(n^2)$}
\label{alg:gsph}
\begin{algorithmic}[1]
\Require $V\!\subset\!\mathbb{R}^2$, $\varepsilon$,
         $T_{\text{bgt}}$, $k_{\text{cand}}$, \textsc{MaxIter}
\Ensure  Tour $\pi$, length $L$
\State \textbf{Phase 1 — Spatial partition}
\State Compute $x_m, y_m$; assign each $v\in V$ to quadrant $Q_k$
\State \textbf{Phase 2 — Intra-quadrant construction}
\ForAll{$Q_k\in\{Q_1,Q_2,Q_3,Q_4\}$}
  \If{$|Q_k|\le 1$}\; $R_k\gets Q_k$
  \Else
    \State $R_k\gets\Call{NearestNeighbor}{Q_k}$
    \For{$t=1,2$}
      \State $R_k\gets\Call{TwoOptPath}{R_k,\textsc{MaxIter}}$
    \EndFor
  \EndIf
\EndFor
\State \textbf{Phase 3 — Inter-quadrant frontier connections}
\ForAll{adjacent pair $(A,B,\text{dir})$}
  \State $(u^*,v^*)\gets\Call{BestFrontierPair}{A,B,\text{dir},\varepsilon}$
  \State Record connection $(u^*,v^*)$
\EndFor
\State \textbf{Phase 4 — Global assembly}
\State $\pi_0\gets\Call{ConcatQuadrants}{\{R_1,R_2,R_3,R_4\}}$
\State \textbf{Phase 5 — Local-search polishing}
\State $\pi\gets\Call{PolishFast}{\pi_0,T_{\text{bgt}},k_{\text{cand}}}$
\State \Return $(\pi,\Call{Length}{\pi})$
\end{algorithmic}
\end{algorithm}

\paragraph{\textsc{NearestNeighbor}.}
Greedy path construction: always select the closest unvisited node.
\textbf{Complexity:} $O(m^2)$, $m=|Q_k|$.

\begin{algorithm}
\caption{\textsc{NearestNeighbor} — Greedy path construction}
\label{alg:nn}
\begin{algorithmic}[1]
\Require Node set $P=\{p_1,\dots,p_m\}$
\Ensure  Open path \textit{route}
\State $\textit{cur}\gets p_1$;\;
       $R\gets P\setminus\{p_1\}$;\;
       $\textit{route}\gets[p_1]$
\While{$R\neq\emptyset$}
  \State $\textit{nxt}\gets\arg\min_{p\in R}d(\textit{cur},p)$
  \State \textit{route}.append(\textit{nxt});\;
         $R\gets R\setminus\{\textit{nxt}\}$;\;
         $\textit{cur}\gets\textit{nxt}$
\EndWhile
\State \Return \textit{route}
\end{algorithmic}
\end{algorithm}

\paragraph{\textsc{TwoOptPath}.}
Applies 2-opt to an open path, reversing segments that reduce total
length. \textbf{Complexity:} $O(m^2)$/iteration.

\begin{algorithm}
\caption{\textsc{TwoOptPath} — 2-opt for open path}
\label{alg:twoopt}
\begin{algorithmic}[1]
\Require Path \textit{route}, \textsc{MaxIter}
\Ensure  Improved path
\State $\textit{imp}\gets\mathrm{true}$;\; $t\gets 0$
\While{$\textit{imp}$ \textbf{and} $t<\textsc{MaxIter}$}
  \State $\textit{imp}\gets\mathrm{false}$
  \For{$i=1$ \textbf{to} $|\textit{route}|-2$}
    \For{$j=i+1$ \textbf{to} $|\textit{route}|-1$}
      \State $\delta\gets d(r_i,r_j)+d(r_{i+1},r_{j+1})
                       -d(r_i,r_{i+1})-d(r_j,r_{j+1})$
      \If{$\delta<0$}
        \State \Call{Reverse}{\textit{route}$[i{+}1{:}j]$};\;
               $\textit{imp}\gets\mathrm{true}$
      \EndIf
    \EndFor
  \EndFor
  \State $t\gets t+1$
\EndWhile
\State \Return \textit{route}
\end{algorithmic}
\end{algorithm}

\paragraph{\textsc{BestFrontierPair}.}
Finds the minimum-cost node pair across adjacent quadrants within
band $\varepsilon$. \textbf{Complexity:} $O(|F_A||F_B|)$.

\begin{algorithm}
\caption{\textsc{BestFrontierPair} — Minimum-cost cross-boundary pair}
\label{alg:frontier}
\begin{algorithmic}[1]
\Require Quadrants $A,B$;\;
         $\texttt{dir}\in\{\text{vert},\text{horiz}\}$;\;
         midpoint $m$;\; tolerance $\varepsilon$
\Ensure  Pair $(u^*,v^*)$
\If{\texttt{dir}=vert}
  \State $F_A\gets\{p\in A:|x_p-m|<\varepsilon\}$;\;
         $F_B\gets\{p\in B:|x_p-m|<\varepsilon\}$
\Else
  \State $F_A\gets\{p\in A:|y_p-m|<\varepsilon\}$;\;
         $F_B\gets\{p\in B:|y_p-m|<\varepsilon\}$
\EndIf
\State $(u^*,v^*)\gets(\bot,\bot)$;\; $d^*\gets+\infty$
\ForAll{$u\in F_A$}
  \ForAll{$v\in F_B$}
    \If{$d(u,v)<d^*$}\;
      $(u^*,v^*,d^*)\gets(u,v,d(u,v))$
    \EndIf
  \EndFor
\EndFor
\State \Return $(u^*,v^*)$
\end{algorithmic}
\end{algorithm}

\paragraph{\textsc{ConcatQuadrants}.}
Concatenates quadrant paths in clockwise order, inverting each
when it reduces the linking distance.
\textbf{Complexity:} $O(n)$.

\begin{algorithm}
\caption{\textsc{ConcatQuadrants} — Quadrant route assembly}
\label{alg:concat}
\begin{algorithmic}[1]
\Require Routes $\{R_1,R_2,R_3,R_4\}$
\Ensure  Concatenated sequence \textit{seq}
\State $\textit{order}\gets[Q_1,Q_2,Q_4,Q_3]$;\;
       $\textit{seq}\gets[\,]$;\;
       $\textit{last}\gets\bot$
\ForAll{$Q_k$ \textbf{in} \textit{order}}
  \State $R\gets$ route of $Q_k$
  \If{$R=\emptyset$}\;\textbf{continue}\EndIf
  \If{$\textit{last}=\bot$}\; $\textit{seq}\gets R$
  \Else
    \If{$d(\textit{last},R[1])\le d(\textit{last},R[\text{end}])$}
      \State Append $R$ to \textit{seq}
    \Else
      \State Append \Call{Reverse}{$R$} to \textit{seq}
    \EndIf
  \EndIf
  \State $\textit{last}\gets\textit{seq}[\text{end}]$
\EndFor
\State \Return \textit{seq}
\end{algorithmic}
\end{algorithm}

\paragraph{\textsc{PolishFast}.}
Global improvement: (i) 2-opt DLB, (ii) Or-opt, (iii) ILS with
double-bridge. Uses candidate sets throughout.
\textbf{Complexity:} $O(n^2)$.

\begin{algorithm}
\caption{\textsc{PolishFast} — Budget-bounded local-search polishing}
\label{alg:polish}
\begin{algorithmic}[1]
\Require Tour $\pi_0$, budget $T_{\text{bgt}}$, $k_{\text{cand}}$
\Ensure  Polished tour
\State $\textit{dl}\gets\Call{Now}{}+T_{\text{bgt}}$;\;
       Build distance matrix $D$ and candidate sets $C$
\State $\pi\gets\Call{2OptDLB}{\pi_0,D,C,\textit{dl}}$
\State $\pi\gets\Call{OrOpt}{\pi,D,C,\textit{dl}}$
\State $\pi\gets\Call{2OptDLB}{\pi,D,C,\textit{dl}}$
\State $\pi^*\gets\pi$;\; $L^*\gets\Call{Length}{\pi,D}$
\For{$t=1$ \textbf{to} \textsc{IlsIters}}
  \If{$\Call{Now}{}>\textit{dl}$}\;\textbf{break}\EndIf
  \State $\pi'\gets\Call{DoubleBridge}{\pi^*}$
  \State $\pi'\gets\Call{2OptDLB}{\pi',D,C,\textit{dl}}$
  \If{$\Call{Length}{\pi',D}<L^*$}\;
    $\pi^*\gets\pi'$;\; $L^*\gets\Call{Length}{\pi',D}$
  \EndIf
\EndFor
\State \Return $\pi^*$
\end{algorithmic}
\end{algorithm}

\paragraph{\textsc{2OptDLB}.}
Accelerated 2-opt using Don't-Look Bits and candidate sets to skip
nodes with no improvement potential.
\textbf{Complexity:} $O(n^2)$.

\begin{algorithm}
\caption{\textsc{2OptDLB} — 2-opt with Don't-Look Bits and candidates}
\label{alg:dlb}
\begin{algorithmic}[1]
\Require Tour $\pi$, distances $D$, candidates $C$,
         deadline \textit{dl}
\Ensure  Locally improved tour
\State Init position array \textit{pos};\;
       $\text{dlb}[v]\gets\mathrm{false}\ \forall v$;\;
       $\textit{imp}\gets\mathrm{true}$
\While{$\textit{imp}$}
  \State $\textit{imp}\gets\mathrm{false}$
  \For{$a=0$ \textbf{to} $|\pi|-1$}
    \If{$\Call{Now}{}>\textit{dl}$}\;\Return $\pi$\EndIf
    \If{$\text{dlb}[\pi_a]$}\;\textbf{continue}\EndIf
    \State $g^*\gets 0$;\; $c^*\gets\bot$
    \ForAll{$c_v\in C[\pi_a]$}
      \State $c\gets\textit{pos}[c_v]$;\;
             $b\gets(a+1)\bmod|\pi|$;\;
             $d\gets(c+1)\bmod|\pi|$
      \State $g\gets D[\pi_a,\pi_b]+D[\pi_c,\pi_d]
                   -D[\pi_a,\pi_c]-D[\pi_b,\pi_d]$
      \If{$g>g^*$}\; $g^*\gets g$;\; $c^*\gets c$\EndIf
    \EndFor
    \If{$g^*>0$}
      \State Apply 2-opt between $b$ and $c^*$;\;
             update \textit{pos};\;
             $\text{dlb}[\pi_a]\gets\mathrm{false}$;\;
             $\textit{imp}\gets\mathrm{true}$
    \Else\; $\text{dlb}[\pi_a]\gets\mathrm{true}$
    \EndIf
  \EndFor
\EndWhile
\State \Return $\pi$
\end{algorithmic}
\end{algorithm}

% ================================================================
\section{Experimental Results}
\label{sec:experiments}
% ================================================================

\paragraph{Setup.}
All runs use the same GSPH-FC configuration across all instances.
Times are single-run wall-clock seconds. GAP$(\%)=100\times
(\text{GSPH}-\text{BKS})/\text{BKS}$; positive values indicate
solutions above the BKS. Instances without a known BKS are marked
\textbf{N/A}. The \texttt{d198} negative GAP is attributed to a
\texttt{CEIL\_2D}/\texttt{EUC\_2D} metric mismatch in the historical
record and is retained for transparency.

\textbf{Small Instances ($n \le 120$)}

\begin{table}
\textbf{Small instances ($n\!\le\!120$).
\textbf{Bold GAP}: optimum or gap $\le\!0.05\%$.}\label{tab:small}
\centering\small
\begin{tabular}{@{}lrrrrr@{}}
\toprule
\textbf{Instance} & $n$ & \textbf{GSPH} & \textbf{BKS}
  & \textbf{GAP (\%)} & \textbf{Time (s)} \\
\midrule
berlin52   &  52 &   7{,}542 &   7{,}542 & \textbf{0.000} & 0.012 \\
eil51      &  51 &       428 &       426 &          0.469 & 0.011 \\
eil76      &  76 &       541 &       538 &          0.558 & 0.034 \\
eil101     & 101 &       634 &       629 &          0.795 & 0.067 \\
eil10      &  10 &       156 &       --- &  \textbf{N/A}  & 0.001 \\
eil20      &  20 &       238 &       --- &  \textbf{N/A}  & 0.002 \\
rd100      & 100 &   7{,}911 &   7{,}910 & \textbf{0.013} & 0.074 \\
rat99      &  99 &   1{,}211 &   1{,}211 & \textbf{0.000} & 0.077 \\
st70       &  70 &       680 &       675 &          0.741 & 0.019 \\
pr76       &  76 & 108{,}159 & 108{,}159 & \textbf{0.000} & 0.040 \\
kroA100    & 100 &  21{,}282 &  21{,}282 & \textbf{0.000} & 0.083 \\
kroB100    & 100 &  22{,}245 &  22{,}141 &          0.470 & 0.034 \\
kroC100    & 100 &  20{,}880 &  20{,}749 &          0.631 & 0.055 \\
kroD100    & 100 &  21{,}446 &  21{,}294 &          0.714 & 0.087 \\
kroE100    & 100 &  22{,}084 &  22{,}068 & \textbf{0.073} & 0.070 \\
lin105     & 105 &  14{,}379 &  14{,}379 & \textbf{0.000} & 0.098 \\
pr107      & 107 &  44{,}765 &  44{,}303 &          1.043 & 0.069 \\
\midrule
\multicolumn{4}{l}{\textit{Avg.\ GAP (excl.\ N/A)}}
  & \textbf{0.313} & \\
\bottomrule
\end{tabular}
\end{table}

\textbf{Medium Instances ($121 \le n \le 300$)}

\begin{table}
\caption{Medium instances ($121\!\le\!n\!\le\!300$).
$^{\dagger}$Negative GAP: likely metric mismatch
(\texttt{CEIL\_2D} vs.\ \texttt{EUC\_2D}).}\label{tab:medium}
\centering\small
\begin{tabular}{@{}lrrrrr@{}}
\toprule
\textbf{Instance} & $n$ & \textbf{GSPH} & \textbf{BKS}
  & \textbf{GAP (\%)} & \textbf{Time (s)} \\
\midrule
bier127        & 127 & 118{,}914 & 118{,}282 &   0.534 & 0.133 \\
ch130          & 130 &   6{,}153 &   6{,}110 &   0.704 & 0.123 \\
ch150          & 150 &   6{,}555 &   6{,}528 &   0.414 & 0.140 \\
pr124          & 124 &  59{,}608 &  59{,}030 &   0.979 & 0.098 \\
pr136          & 136 & 100{,}193 &  96{,}772 &   3.535 & 0.140 \\
pr144          & 144 &  65{,}960 &  58{,}537 &  12.684 & 0.167 \\
pr152          & 152 &  75{,}203 &  73{,}682 &   2.064 & 0.147 \\
ts225          & 225 & 128{,}485 & 126{,}643 &   1.454 & 0.296 \\
tsp225         & 225 &   4{,}010 &   3{,}916 &   2.400 & 0.379 \\
a280           & 280 &   2{,}652 &   2{,}579 &   2.831 & 0.488 \\
kroA150        & 150 &  27{,}297 &  26{,}524 &   2.914 & 0.237 \\
kroB150        & 150 &  27{,}615 &  26{,}130 &   5.683 & 0.163 \\
pr226          & 226 &  85{,}729 &  80{,}369 &   6.667 & 0.350 \\
pr264          & 264 &  50{,}974 &  49{,}135 &   3.740 & 0.037 \\
pr299          & 299 &  48{,}749 &  48{,}191 &   1.158 & 0.953 \\
kroA200        & 200 &  29{,}658 &  29{,}368 &   0.987 & 0.370 \\
kroB200        & 200 &  30{,}097 &  29{,}437 &   2.241 & 0.410 \\
gil262         & 262 &   2{,}669 &   2{,}378 &  12.240 & 0.849 \\
d198$^\dagger$ & 198 &  13{,}738 &  15{,}780 &$-$12.940 & 0.491 \\
u159           & 159 &  42{,}080 &  42{,}080 & \textbf{0.000} & 0.177 \\
\midrule
\multicolumn{4}{l}{\textit{Avg.\ GAP (excl.\ $^\dagger$)}}
  & \textbf{3.381} & \\
\bottomrule
\end{tabular}
\end{table}

\textbf{Large Instances ($301 \le n \le 2000$)}

\begin{table}
\caption{Large instances ($301\!\le\!n\!\le\!2000$).}\label{tab:large}
\centering\small
\begin{tabular}{@{}lrrrr@{}}
\toprule
\textbf{Instance} & $n$ & \textbf{GAP (\%)} & \textbf{Time (s)} \\
\midrule
lin318    &     318 &  1.98 &    1.4 \\
linhp318  &     318 &  4.36 &    1.1 \\
fl417     &     417 &  6.32 &    2.2 \\
pcb442    &     442 &  9.95 &    2.6 \\
rd400     &     400 &  8.34 &    2.6 \\
pr439     &     439 &  9.31 &    3.8 \\
p654      &     654 &  9.26 &    9.8 \\
rat575    &     575 & 10.16 &    4.2 \\
u574      &     574 &  7.25 &    6.7 \\
rat783    &     783 &  9.39 &    9.5 \\
u724      &     724 &  9.56 &   10.0 \\
d493      &     493 &  4.05 &    3.1 \\
d657      &     657 &  7.74 &    7.5 \\
\midrule
pr1002    & 1{,}002 &  9.08 &   20.3 \\
u1060     & 1{,}060 &  8.42 &   28.8 \\
pcb1173   & 1{,}173 & 11.00 &   25.4 \\
nrw1379   & 1{,}379 &  7.83 &   36.0 \\
rl1304    & 1{,}304 & 11.85 &   40.8 \\
rl1323    & 1{,}323 & 11.08 &   30.1 \\
rl1889    & 1{,}889 & 10.25 &   92.7 \\
d1291     & 1{,}291 &  5.83 &   26.2 \\
d1655     & 1{,}655 &  6.60 &   61.6 \\
fl1400    & 1{,}400 &  4.05 &   39.7 \\
fl1577    & 1{,}577 &  9.68 &   73.0 \\
u1432     & 1{,}432 &  6.64 &   35.4 \\
u1817     & 1{,}817 & 10.67 &   63.5 \\
vm1084    & 1{,}084 & 10.10 &   22.6 \\
vm1748    & 1{,}748 &  8.95 &   94.7 \\
\midrule
\multicolumn{2}{l}{\textit{Avg.\ GAP}} & \textbf{8.14} & \\
\bottomrule
\end{tabular}
\end{table}

\textbf{Very Large Instances ($n > 2000$)}

\begin{table}
\caption{Very large instances ($n\!>\!2000$).}\label{tab:xlarge}
\centering\small
\begin{tabular}{@{}lrrrr@{}}
\toprule
\textbf{Instance} & $n$ & \textbf{GAP (\%)} & \textbf{Time (s)} \\
\midrule
pr2392  & 2{,}392 &  7.41 &     126.1 \\
u2152   & 2{,}152 &  9.77 &     114.8 \\
u2319   & 2{,}319 &  2.59 &     125.0 \\
pcb3038 & 3{,}038 &  7.91 &     319.8 \\
fl3795  & 3{,}795 &  3.63 &     382.2 \\
fnl4461 & 4{,}461 &  8.05 &     787.4 \\
rl5915  & 5{,}915 &  7.66 & 1{,}510.0 \\
rl5934  & 5{,}934 &  8.43 & 1{,}366.6 \\
\midrule
\multicolumn{2}{l}{\textit{Avg.\ GAP}} & \textbf{6.93} & \\
\bottomrule
\end{tabular}
\end{table}

\textbf{Visual Consistency Analysis}

To assess structural stability, GSPH-FC was executed 10 times
independently on \texttt{berlin52}, \texttt{eil76}, and
\texttt{ts225}. Figures~\ref{fig:berlin52}--\ref{fig:ts225} display
the resulting tours, enabling qualitative comparison of routing
patterns and inter-run variability. This analysis allows observing
whether the method converges to similar tour topologies or exhibits
high geometric dispersion across runs.

\begin{figure}
\centering
\includegraphics[width=\textwidth]%
  {Resultados/grafico_berlin52_comparison.jpg}
\caption{Ten independent GSPH-FC runs on \texttt{berlin52} ($n=52$).
Strong structural consistency on a small, well-structured instance.}
\label{fig:berlin52}
\end{figure}

\begin{figure}
\centering
\includegraphics[width=\textwidth]%
  {Resultados/grafico_eil76_comparison.jpg}
\caption{Ten independent runs on \texttt{eil76} ($n=76$). Moderate
geometric variability; quadrant boundary patterns are recognizable
across all runs.}
\label{fig:eil76}
\end{figure}

\begin{figure}
\centering
\includegraphics[width=\textwidth]%
  {Resultados/grafico_ts225_comparison.jpg}
\caption{Ten independent runs on \texttt{ts225} ($n=225$). Structural
robustness is clearly visible; local differences arise mainly near
quadrant boundaries.}
\label{fig:ts225}
\end{figure}

\subsection{Peak Memory Analysis}
\label{sec:memory}

Peak memory was profiled with \texttt{tracemalloc} and
\texttt{pympler} on the same three instances used for visual
analysis. Table~\ref{tab:memory} isolates the GSPH-specific
heap allocations at \texttt{GSPH\_BASE.py} lines 58 and 185
(route and candidate structures), and reports the total process
peak which includes Python runtime and profiler overhead
($\approx$1.65\,MB constant across all instances).

\begin{table}
\caption{Peak memory profile across five representative instances.
\textit{GSPH alloc.}\ sums all algorithm-specific hotspots
in \texttt{GSPH\_BASE.py}; process peak includes Python runtime
and profiler overhead ($\approx$1.65\,MB constant).}\label{tab:memory}
\centering\small
\begin{tabular}{@{}lrrrrr@{}}
\toprule
\textbf{Instance} & $n$
  & \textbf{Tuples (MB)} & \textbf{float (KB)}
  & \textbf{GSPH alloc.\ (KiB)} & \textbf{Peak (MB)} \\
\midrule
berlin52 &  52 & 2.74 & 1.27 & $23.1{+}2.9 \approx 26.0$ & $\approx 4.82$ \\
eil76    &  76 & 2.73 & 1.27 & $20.4{+}3.9 \approx 24.3$ & $\approx 4.81$ \\
pr152    & 152 & 2.74 & 1.83 & $20.7{+}3.9 \approx 24.6$ & $\approx 4.82$ \\
ts225    & 225 & 2.77 & 4.64 & $27.6{+}3.9 \approx 31.5$ & $\approx 4.85$ \\
\bottomrule
\end{tabular}
\end{table}

Four key observations emerge from the memory profile.
%
\textbf{(1) Near-constant total footprint.}
The process peak spans only 4.81--4.85\,MB across $n=76$ to $n=225$,
a spread of $\approx$40\,KB for a $3\times$ increase in instance
size. The dominant consumer is the Python runtime itself (tuples,
lists, strings) at a constant $\approx$4.63\,MB, confirming that
GSPH-FC introduces virtually no additional heap pressure as $n$ grows.
%
\textbf{(2) Sub-linear GSPH-specific growth.}
Algorithm-specific allocations range from 24.3\,KiB ($n=76$) to
31.5\,KiB ($n=225$), a $1.30\times$ factor for a $3\times$ increase
in $n$. This sub-linear scaling is consistent with the quadrant
decomposition: each sub-problem holds $\approx n/4$ nodes, so local
route and candidate-set structures grow as $O(n)$ rather than $O(n^2)$.
%
\textbf{(3) Absence of a full distance matrix.}
Float objects stay below 5\,KiB across all instances. A full
\texttt{float64} distance matrix for $n=225$ would cost
$225^2\times 8\approx 405$\,KiB — a $80\times$ overhead that
never materializes, confirming that GSPH-FC computes distances
on-the-fly and retains only sparse candidate-set entries.
%
\textbf{(4) Instance-independent list and string overhead.}
Both \texttt{list} (1.22\,MB) and \texttt{str} ($\approx$687\,KB)
counts are identical across all four instances, confirming they
belong to the Python interpreter and profiler infrastructure rather
than the algorithm. This property guarantees that switching to a
leaner runtime (e.g., PyPy or a compiled extension) would reduce
the baseline overhead without changing the algorithm's $O(n)$
memory profile.


% ================================================================
\section{Critical Analysis}
\label{sec:analysis}
% ================================================================

\paragraph{Strengths.}
\begin{itemize}
  \item \textit{Effective spatial locality.} Quadrant decomposition
    confines 2-opt to compact sub-spaces, reducing crossing edges
    and yielding interpretable tours.
  \item \textit{Natural parallelism.} Each $Q_k$ is solved
    independently (NN + 2-opt), enabling 4-thread execution with
    near-linear speedup.
  \item \textit{Reduced constant factor.} Working on sub-tours of
    $\approx n/4$ nodes reduces the practical 2-opt cost to
    $\sim O(n^2/4)$.
  \item \textit{Simplicity and traceability.} Clear frontier rules
    simplify debugging, profiling, and integration in production
    pipelines.
  \item \textit{Composability.} The polishing phase (DLB +
    candidate sets + ILS) integrates seamlessly without redesign.
  \item \textit{Memory efficiency.} Near-constant peak footprint
    ($\approx$4.75\,MB) makes the method suitable for
    memory-constrained deployments.
\end{itemize}

\paragraph{Weaknesses.}
\begin{itemize}
  \item \textit{Rigid axis-aligned cuts.} Clusters straddling
    $x_m$ or $y_m$ may be split, inducing long linking edges.
  \item \textit{Distribution sensitivity.} Highly imbalanced
    quadrants degrade both time and assembly quality.
  \item \textit{Fixed concatenation order.} A static clockwise
    sweep can produce zig-zag patterns on non-uniform instances.
  \item \textit{Metric dependence.} Switching between
    \texttt{EUC\_2D} and \texttt{CEIL\_2D} may produce atypical
    GAPs if BKS records are inconsistent.
\end{itemize}

\paragraph{Opportunities.}
\begin{itemize}
  \item \textit{Boundary hybridization.} Applying 3-opt or a
    focused LK pass only on frontier nodes corrects long links at
    low cost.
  \item \textit{Recursive QML-GSPH.} A quadtree decomposition with
    adaptive stopping threshold $n_{\min}$ extends the hierarchical
    idea to multiple levels.
  \item \textit{Delaunay/MST candidate sets.} Geometry-guided
    $k$-NN structures accelerate 2-opt and Or-opt without
    sacrificing improvement quality.
  \item \textit{Adaptive ordering.} Deciding the concatenation
    sequence by density centroid or PCA reduces zig-zag artifacts.
\end{itemize}

\paragraph{Threats.}
\begin{itemize}
  \item \textit{Non-Euclidean metrics.} Asymmetric or road-network
    instances violate the geometric assumptions of GSPH-FC.
  \item \textit{Adversarial distributions.} Bimodal or diagonal
    clusters increase frontier connection costs significantly.
  \item \textit{BKS inconsistencies.} Metric or rounding
    discrepancies in the reference library can skew GAP reporting.
\end{itemize}

\paragraph{CPU behavior.}
CPU usage tends to increase during the polishing phase
(2-opt/Or-opt/ILS) as the algorithm explores harder neighborhoods
after large improvements stabilize. This is mitigated by: limiting
$k\in[20,40]$ in candidate sets with aggressive DLB; per-phase
time quotas with early stopping; restricting 3-opt/LK to frontier
windows; rebalancing via PCA rotation when quadrant sizes are
imbalanced; precomputing $D$ in \texttt{float32} and vectorizing
with Numba/Cython/OpenMP.

% ================================================================
\section{Conclusions}
\label{sec:conclusions}
% ================================================================

We presented \textbf{GSPH-FC}, a hierarchical Euclidean TSP
heuristic combining bounding-box quadrant decomposition, Nearest
Neighbour + 2-opt per quadrant, and minimum-cost frontier
connectivity for global assembly. On 55 TSPLIB instances
($n$ up to 5\,934), the method achieves a typical GAP of 3--12\%
over the BKS with sub-second times for $n\!\le\!300$, under 100\,s
for $n\!\le\!1{,}900$, and natural 4-thread parallelism that
reduces the effective 2-opt constant factor by $\approx\!4\times$.
Memory profiling confirms a near-constant footprint ($\approx$4.75\,MB)
dominated by Python runtime overhead, with GSPH-specific allocations
growing sub-linearly at just 31.5\,KiB for $n=225$.

The geometric design produces interpretable, low-crossing tours and
supports straightforward integration in logistics deployments.
One atypical case (\texttt{d198}, negative GAP) highlights the
critical importance of consistent metric conventions across
experimental pipelines.

Future work includes: (i)~recursive \textbf{QML-GSPH} with
density-adaptive $n_{\min}$; (ii)~focused 3-opt/LK on frontier
edges; (iii)~PCA- or centroid-guided adaptive concatenation;
and (iv)~extension to asymmetric TSP, road-network VRP, and
VRPTW where spatial decomposition can serve as a geometric
pre-solver.

\begin{credits}
\subsubsection{\ackname}
The authors thank the faculty of the Heuristics course at
Universidad Andres Bello for their guidance and material support
during the development of this work.
\end{credits}

\bibliographystyle{splncs04}
\bibliography{bibliografia}

\end{document}